\uuid{zfWy}
\titre{Équations complexes}
\auteur{bodin ; hanani ; mzari}
\organisation{exo7}
\licence{CC BY-NC-SA 4.0}
\datecreate{2020-01-20}
\niveau{L1}
\module{Algèbre}
\chapitre{Nombres complexes}
\sousChapitre{Racine carrée, équation du second degré}
\difficulte{2}
\theme{Équations, équations complexes}
\exo7qcmid{lille-1.nombres-complexes.020}

\contenu{
\begin{question}
\qclassification{104.02,104.03,104.04}
\qtags{motcle=équations complexes}
On considère l'équation : $(E) : \, z^2-2iz-1-i=0$, $z\in \Cc$.   Quelles sont les assertions vraies ?
\begin{answers}
    \bad{Le discriminant de l'équation est : $\Delta = 8+4i$.}

    \good{Le discriminant de l'équation est : $\Delta = 4i$.}

    \bad{les solutions de $(E)$ sont :  $z_1=\frac{\sqrt 2+ (1+\sqrt 2)i}{2}$ et $z_2=\frac{\sqrt 2+ (1-\sqrt 2)i}{2}$.}

    \good{les solutions de $(E)$ sont : $z_1=\frac{\sqrt 2+ (2+\sqrt 2)i}{2}$ et $z_2=\frac{-\sqrt 2+ (2-\sqrt 2)i}{2}$.}
\end{answers}
\begin{explanations}
Utiliser la méthode de résolution d'une équation du second degré.
\end{explanations}
\end{question}
}
